\documentclass[11pt]{article}
\usepackage[margin=1in]{geometry}
\usepackage{times}
\usepackage{hyperref}
\hypersetup{colorlinks=true, linkcolor=black, urlcolor=blue, citecolor=black}
\title{The Device That Was Smashed With a Hammer: T. Henry Moray's Radiant Energy Device}
\author{Flyxion \\ \small Independent Researcher}
\date{}
\begin{document}
\maketitle

\begin{abstract}
T. Henry Moray demonstrated, from the 1920s through the 1930s, a "Radiant Energy Device" claimed to draw substantial usable electrical power from ambient cosmic and atmospheric radiation using a specially prepared germanium-based semiconductor detector, a claim that ended without independent resolution when Moray's prototype was destroyed, by Moray's own account through sabotage, in an incident for which no independent corroborating evidence beyond Moray's own testimony has ever surfaced. This paper argues that the specific germanium detector technology Moray worked with, while a genuine and historically important precursor to later semiconductor development, has no established physical basis for extracting the claimed power levels from ambient radiation, that Moray's own inconsistent and shifting technical descriptions across different demonstrations undermine rather than support a single, well-specified underlying mechanism, and that the destroyed-prototype narrative, whatever its truth, has functioned since to place the central evidentiary artifact permanently outside the reach of independent examination in a manner directly comparable to other undemonstrated devices surveyed elsewhere in this series.
\end{abstract}

\section{Introduction}

T. Henry Moray, a Utah engineer, demonstrated a device from the 1920s onward claimed to detect and convert ambient cosmic ray and radio frequency background radiation into usable electrical power at levels sufficient to run substantial electrical loads, using an arrangement centered on a germanium crystal detector, radium-treated components, and a resonant antenna and tuning circuit, presented to investors and occasional independent witnesses over roughly two decades, culminating in a well documented 1939 incident in which Moray's prototype was destroyed, reportedly smashed with a hammer by an unidentified intruder, an event Moray attributed to sabotage by unnamed competing commercial interests threatened by the technology, after which no comparably capable device was ever again demonstrated before Moray's death in 1974.

\section{Why the Claimed Power Levels Do Not Follow From the Detector Technology Described}

Germanium crystal detectors, a genuine and historically significant early semiconductor technology that Moray's work anticipated in some respects, function by rectifying weak radio frequency signals, converting an oscillating input into a small direct current output, a real and useful property exploited in early crystal radio receivers, but one that converts and rectifies energy already present in an incoming signal rather than generating new usable power from ambient background radiation at the sustained kilowatt-range output levels some of Moray's demonstrations claimed; ambient cosmic ray and radio background radiation carries an extremely low power density, and no rectification mechanism, however cleverly constructed, can extract more power from a signal than the signal itself carries, meaning the claimed output magnitude requires either an ambient radiation power density many orders beyond what is independently measured to exist in the relevant frequency ranges, or an additional power source beyond simple rectification of ambient background signal that Moray's technical descriptions do not specify.

\section{The Shifting Technical Description Across Demonstrations}

Moray's own accounts of the device's operating principle varied across different demonstrations and written descriptions over the decades he promoted it, at times emphasizing cosmic ray detection, at other times radio frequency reception, and at still other points invoking more general and less specific "radiant energy" language without a single, stable, consistently maintained technical mechanism across his own presentations, a pattern that makes independent technical evaluation of what, precisely, is being claimed considerably more difficult than for a device with a single, stable, published technical description, and that is more consistent with an inventor working from an evolving, not fully specified intuition than with a single well-understood underlying physical mechanism being consistently and accurately described across two decades of promotion.

\section{The Destroyed Prototype and the Sabotage Narrative}

The 1939 destruction of Moray's prototype rests entirely on Moray's own account, with no independent police report, witness testimony beyond Moray's own household, or corroborating physical evidence, such as recovered tool marks consistent with forced entry, that has been independently verified by researchers examining the historical record; regardless of whether the specific incident occurred as described, its effect has been to permanently remove the central physical evidence for the device's claimed function from any possibility of independent post-hoc examination, a structural outcome, evidence permanently unavailable at the exact moment independent verification would matter most, that recurs across multiple devices surveyed in this series regardless of whether sabotage, accident, or simple non-preservation is the actual explanation in any individual case.

\section{The Entropic Gravity Framing}

Moray's radiant energy claim concerns power extraction from ambient background radiation rather than gravitational coupling directly, and is included in this series as an early historical comparative case in the broader "ambient energy harvesting at implausible power density" genre alongside the zero-point energy entries addressed elsewhere here; the core objection, that no described rectification mechanism can extract more power from a signal than the signal itself carries, applies independent of any gravitational framework.

\section{Objections}

Supporters point to Moray's continued and consistent personal conviction across decades, along with a small number of contemporary witness accounts describing having seen the device power incandescent bulbs and other loads, as evidence the underlying demonstrations were not simply fraudulent. Personal conviction sustained across decades and informal witness testimony describing an impressive demonstration are consistent with, and do not distinguish between, a device that genuinely worked as claimed and a device whose actual power source, whether a concealed battery, mains connection, or another conventional source not disclosed to observers, was not identified by witnesses lacking the specific technical background or opportunity to independently inspect the device's full wiring and power path during the demonstration, a distinction the available historical testimony does not resolve.

\section{Conclusion}

The germanium detector technology at the center of Moray's claim provides a mechanism for rectifying, not generating, electrical power, and cannot on its own account for the sustained high power outputs some demonstrations claimed given the low power density of ambient background radiation, Moray's own technical description of the device's operating principle shifted across decades of promotion rather than remaining stable and precisely specified, and the prototype's destruction, resting entirely on Moray's own uncorroborated account, removed the central evidentiary artifact from independent examination at precisely the point such examination would have been most valuable.

\end{document}
